% =====================================================================
%  centro_oftalmologico_talca_ficticio.tex
%  Información general, políticas, catálogo de exámenes, reglas de
%  disponibilidad, precios y protocolos de atención.
%  Compilar con:  pdflatex centro_oftalmologico_talca_ficticio.tex
%  (dos pasadas, para resolver la tabla de contenido)
% =====================================================================
\documentclass[11pt,a4paper]{article}

% ------------------------------------------------- extraccion de texto
\usepackage{cmap}
\usepackage[T1]{fontenc}
\usepackage[utf8]{inputenc}
\usepackage{lmodern}
\usepackage{microtype}
\DisableLigatures{encoding = *, family = *}
\input{glyphtounicode}
\pdfgentounicode=1

% ----------------------------------------------------------------- paquetes
\usepackage[es-nodecimaldot,es-noshorthands,spanish]{babel}
\usepackage[margin=2.2cm]{geometry}
\usepackage[table]{xcolor}
\usepackage{booktabs}
\usepackage{tabularx}
\usepackage{longtable}
\usepackage{enumitem}
\usepackage{fancyhdr}
\usepackage[hidelinks]{hyperref}

\setcounter{secnumdepth}{2}
\setcounter{tocdepth}{2}
\setlength{\parindent}{0pt}
\setlength{\parskip}{4pt}

% ----------------------------------------------------------------- colores
\definecolor{colverde}{HTML}{1B5E20}
\definecolor{fondoverde}{HTML}{E8F5E9}
\definecolor{colrojo}{HTML}{B71C1C}
\definecolor{fondorojo}{HTML}{FDECEA}
\definecolor{colazul}{HTML}{1A237E}
\definecolor{fondoazul}{HTML}{E8EAF6}
\definecolor{fondotabla}{HTML}{ECEFF1}
\definecolor{headcol}{HTML}{37474F}

% ----------------------------------------------------------------- cajas
\newsavebox{\boxdatos}
\newsavebox{\boxdest}
\newsavebox{\boxnota}

\newenvironment{CajaDatos}
  {\par\medskip\begin{lrbox}{\boxdatos}%
   \begin{minipage}{0.92\textwidth}\small}
  {\end{minipage}\end{lrbox}%
   \noindent\setlength{\fboxsep}{8pt}\setlength{\fboxrule}{0.8pt}%
   \fcolorbox{colverde}{fondoverde}{\usebox{\boxdatos}}\par\medskip}

\newenvironment{CajaDestacada}
  {\par\medskip\begin{lrbox}{\boxdest}%
   \begin{minipage}{0.92\textwidth}\small}
  {\end{minipage}\end{lrbox}%
   \noindent\setlength{\fboxsep}{8pt}\setlength{\fboxrule}{0.8pt}%
   \fcolorbox{colrojo}{fondorojo}{\usebox{\boxdest}}\par\medskip}

\newenvironment{CajaNota}
  {\par\medskip\begin{lrbox}{\boxnota}%
   \begin{minipage}{0.92\textwidth}\small}
  {\end{minipage}\end{lrbox}%
   \noindent\setlength{\fboxsep}{8pt}\setlength{\fboxrule}{0.8pt}%
   \fcolorbox{colazul}{fondoazul}{\usebox{\boxnota}}\par\medskip}

\newcommand{\boxtitle}[1]{\textbf{\color{headcol}#1}\par\vspace{2pt}}

% -------------------------------------------------- atajos de contenido
\newcommand{\pesos}[1]{\mbox{\$ #1}}

\newlist{campos}{description}{1}
\setlist[campos]{style=nextline,font=\normalfont\bfseries,
                  itemsep=3pt,topsep=4pt}

\newlist{hitos}{itemize}{1}
\setlist[hitos]{label=\textbullet,leftmargin=1.4em,itemsep=2pt,topsep=3pt}

\newlist{planos}{itemize}{1}
\setlist[planos]{label=\textbullet,leftmargin=1.4em,itemsep=2pt,topsep=3pt}

% ------------------------------------------------------ encabezado y pie
\pagestyle{fancy}
\fancyhf{}
\renewcommand{\headrulewidth}{0.4pt}
\fancyhead[L]{\footnotesize\color{headcol}Centro Oftalmológico Talca}
\fancyhead[R]{\footnotesize\color{headcol}%
  Información, políticas y catálogo de exámenes}
\fancyfoot[C]{\footnotesize\color{headcol}\thepage}

% =====================================================================
\begin{document}

% ----------------------------------------------------------------- portada
\begin{titlepage}
\centering

\vspace*{3.2cm}

{\Large\bfseries\color{headcol} Centro Oftalmológico}\\[0.4em]
{\Huge\bfseries\color{headcol} Talca, Chile}

\vspace{1.4cm}
{\large\color{headcol} Información del centro, políticas, catálogo de
exámenes,\\ disponibilidad y precios}

\vfill
{\small\color{headcol} Actualizado el \today}
\end{titlepage}

% ------------------------------------------------------- tabla de contenido
\tableofcontents
\newpage

% =====================================================================
\section{Información general del centro}

\subsection{Identificación}
\label{sec:ident}

\begin{CajaDatos}
\begin{campos}
\item[Nombre] Centro Oftalmológico Talca.
\item[Dirección] 3 Oriente 1608, Talca, Región del Maule, Chile.
\item[Teléfono] +56 71 223 2252.
\item[Nombre comercial] Centro Oftalmológico Talca.
\item[Razón social] Centro Oftalmológico Talca SpA.
\item[RUT] 77.000.000-0.
\item[Registro de prestador] PS-0000. Código de registro interno.
\end{campos}
\end{CajaDatos}

\subsection{Canales de contacto}

\begin{CajaDatos}
\begin{center}
\begin{tabularx}{\textwidth}{@{}p{3.2cm}p{5.0cm}X@{}}
\toprule
\rowcolor{fondotabla}\textbf{Canal} & \textbf{Dato} & \textbf{Uso} \\
\midrule
Teléfono & +56 71 223 2252 &
  Llamada de recepción. Entrega información de valores y horarios. \\
WhatsApp & +56 9 6123 7788 &
  Mensajes escritos y envío de la orden médica. Se responde en horario de
  atención. \\
Correo & \url{cotizacion@oftalca-centro-demo.example} &
  Envío de ordenes médicas y solicitud de informes. \\
Formulario web & \url{cotizacion.oftalca-centro-demo.example} &
  Solicitud de hora sin llamada. \\
\bottomrule
\end{tabularx}
\end{center}
\end{CajaDatos}

\subsection{Horarios y días de atención}
\label{sec:horarios}

\begin{CajaDatos}
\begin{center}
\begin{tabularx}{\textwidth}{@{}p{3.4cm}p{4.6cm}X@{}}
\toprule
\rowcolor{fondotabla}\textbf{Bloque} & \textbf{Días} & \textbf{Horas} \\
\midrule
Atención general & Lunes a viernes & 08:30 a 13:00 \\
Tarde & Martes, jueves y viernes & 15:00 a 19:00 \\
Atención sábado & Sábado & 09:00 a 13:00 \\
Cerrado & Domingo y feriados & Sin atención \\
\bottomrule
\end{tabularx}
\end{center}

El bloque de la tarde existe solo en parte de la semana: los días sin tarde
(lunes, miércoles y sábado) no tienen horas por la tarde.
\end{CajaDatos}

\subsection{Horarios diferenciados por examen}

\begin{CajaDatos}
\begin{center}
\begin{tabularx}{\textwidth}{@{}p{2.9cm}p{4.0cm}X@{}}
\toprule
\rowcolor{fondotabla}\textbf{Código} & \textbf{Examen} &
\textbf{Bloque reservado} \\
\midrule
OFT-CV-001 & Curvimetría & Lunes a sábado, en horario general. \\
OFT-GLU-002 & Medición de glaucoma & Solo martes, jueves y viernes, bloques
  de mañana. \\
OFT-FDO-003 & Fondo de ojo & Lunes a viernes en mañana, más miércoles por la
  tarde. \\
\bottomrule
\end{tabularx}
\end{center}
\end{CajaDatos}

\subsection{Canales disponibles para consultas}

\begin{CajaDatos}
\begin{hitos}
\item \textbf{Teléfono:} canal principal. Entrega valores y horarios en la
  misma llamada en la mayoría de los casos.
\item \textbf{WhatsApp:} canal secundario. Entrega los mismos valores, pero
  la disponibilidad solo se informa si el mensaje llega antes de las 17:00.
\item \textbf{Correo:} canal para enviar ordenes médicas. No se responde
  sobre valores por correo.
\item \textbf{Formulario web:} canal para solicitud de hora. La confirmación
  llega en un mensaje al WhatsApp registrado.
\end{hitos}
\end{CajaDatos}

\subsection{Modalidades de atención}

\begin{CajaDatos}
\begin{center}
\begin{tabularx}{\textwidth}{@{}p{4.0cm}X@{}}
\toprule
\rowcolor{fondotabla}\textbf{Modalidad} & \textbf{Descripción} \\
\midrule
Presencial & Toda la atención se realiza en el domicilio del centro, en la
  dirección indicada en la sección \ref{sec:ident}. Es la uúnica modalidad
  que realizan los tres exámenes del catálogo. \\
Telemedicina & Solo para la toma de antecedentes clínicos previos a
  OFT-GLU-002, en un bloque de 10 minutos. No reemplaza el examen
  presencial. \\
Entrega de resultados & El informe se entrega en mano al terminar, o por
  correo electrónico a petición del paciente. \\
\bottomrule
\end{tabularx}
\end{center}
\end{CajaDatos}

\subsection{Anticipación recomendada para solicitar una hora}

\begin{CajaDatos}
\begin{center}
\begin{tabularx}{\textwidth}{@{}p{3.8cm}p{3.4cm}X@{}}
\toprule
\rowcolor{fondotabla}\textbf{Examen} & \textbf{Anticipación} &
\textbf{Con menos anticipación} \\
\midrule
Curvimetría & 24 horas & Se ofrece el mismo día si queda espacio antes de
  las 11:00. \\
Medición de glaucoma & 5 días hábiles & El centro informa que no puede
  garantizar hora y pide volver a llamar. \\
Fondo de ojo & 3 días hábiles & Se ofrece hora en el mismo bloque de la
  tarde, si el paciente puede permanecer cuatro horas con visión borrosa. \\
\bottomrule
\end{tabularx}
\end{center}
\end{CajaDatos}

\subsection{Condiciones generales de atención}

\begin{CajaDestacada}
\begin{planos}
\item El paciente debe presentarse 15 minutos antes de la hora.
\item La hora se respeta según el orden de llegada a recepción, no según el
  orden de la agenda. Si el paciente llega tarde, entra en la cola del bloque.
\item No se realiza ningún examen sin el registro previo en recepción.
\item El centro no atiende urgencias oftalmológicas. Ante un síntoma reciente
  de visión borrosa, dolor ocular o pérdida de visión, el paciente debe
  acudir a un servicio de urgencias.
\item El personal del centro se identifica con nombre y cargo al iniciar la
  atención.
\end{planos}
\end{CajaDestacada}

\subsection{Política de cancelación}

\begin{CajaDestacada}
\begin{center}
\begin{tabularx}{\textwidth}{@{}p{4.6cm}X@{}}
\toprule
\rowcolor{fondorojo}\textbf{Momento de la cancelación} & \textbf{Efecto} \\
\midrule
Con más de 48 horas & Sin costo. Se devuelve la hora al sistema. \\
Entre 24 y 48 horas & Sin costo, pero se registra una cancelación en el
  historial del paciente. \\
Con menos de 24 horas & Se cobra el 50 por ciento del valor del examen,
  porque el bloque queda bloqueado para el resto del día. \\
No se presenta & Se cobra el 100 por ciento del valor del examen. \\
\bottomrule
\end{tabularx}
\end{center}
\end{CajaDestacada}

\subsection{Política de modificación de hora}

\begin{CajaDestacada}
\begin{planos}
\item Una modificación solicitada con más de 24 horas de anticipación no
  tiene costo.
\item Una modificación solicitada con menos de 24 horas se rige por la misma
  regla que la cancelación tardía: 50 por ciento del valor.
\item El paciente puede modificar su hora \textbf{una sola vez} sin costo.
\item Si el paciente modifica dos veces, el centro exige confirmar por
  teléfono antes de aceptar la tercera modificación.
\item El centro nunca llama para ofrecer una hora distinta. Si la hora
  original cae, el paciente debe volver a consultar.
\end{planos}
\end{CajaDestacada}

\subsection{Condiciones para pacientes que llegan atrasados}

\begin{CajaDestacada}
\boxtitle{Tolerancia máxima: 15 minutos}
\begin{center}
\begin{tabularx}{\textwidth}{@{}p{4.2cm}X@{}}
\toprule
\rowcolor{fondorojo}\textbf{Atraso} & \textbf{Efecto} \\
\midrule
Hasta 15 minutos & Se atiende normalmente, con el tiempo restante del bloque. \\
Entre 15 y 30 minutos & El examen se realiza, pero puede quedar incompleto:
  en el caso de OFT-GLU-002 el paciente sale del bloque sin haber hecho la
  campimetría y debe reagendar solo esa parte. \\
Mas de 30 minutos & No se realiza el examen. Se aplica la política de
  no presentación: cobro del 100 por ciento. \\
\bottomrule
\end{tabularx}
\end{center}
\end{CajaDestacada}

\subsection{Documentación requerida}

\begin{CajaDestacada}
\begin{planos}
\item \textbf{Toda atención:} documento de identidad del paciente o de su
  responsable, para el registro en recepción.
\item \textbf{Todo examen:} cédula vigente o pasaporte vigente. No se acepta
  otro documento.
\item \textbf{Examen con orden médica:} la orden médica original, en papel,
  con firma y timbre del médico. No se acepta una foto en el celular.
\item \textbf{Paciente menor de 18 años:} cédula del adulto responsable y
  autorización firmada de ese adulto.
\end{planos}
\end{CajaDestacada}

\subsection{Ordenes médicas}

\begin{CajaDestacada}
\begin{center}
\begin{tabularx}{\textwidth}{@{}p{4.4cm}X@{}}
\toprule
\rowcolor{fondorojo}\textbf{Regla} & \textbf{Detalle} \\
\midrule
Vigencia & La orden médica tiene una vigencia de 90 días desde su fecha de
  emisión. \\
Datos mínimos & Debe indicar el examen a realizar, la fecha y el médico
  firmante. No se acepta una orden de un examen distinto al que el paciente
  quiere. \\
Quién la emite & Puede ser cualquier médico habilitado, no necesariamente
  oftalmólogo. \\
Sin orden & OFT-CV-001 no la exige. OFT-GLU-002 la exige siempre.
  OFT-FDO-003 la exige solo si el paciente quiere usar Fonasa o convenio. \\
Edad del paciente & Si el paciente tiene 70 años o más, la orden puede ser
  emitida por un médico no especialista. \\
Envío previo & La orden puede enviarse por correo hasta 24 horas antes de la
  hora, para evitar la fila. \\
\bottomrule
\end{tabularx}
\end{center}
\end{CajaDestacada}

\subsection{Entrega de resultados}

\begin{CajaDestacada}
\begin{planos}
\item Ningun resultado se entrega por teléfono ni por WhatsApp.
\item El informe se entrega en mano al terminar el examen, en un sobre
  sellado con el nombre del paciente.
\item Si el paciente no retira el informe en el momento, queda disponible en
  recepción por 60 días.
\item Pasados los 60 días, el informe se elimina y no se puede recuperar:
  hay que pedir una nueva realización del examen.
\item Las imagenes digitalizadas se entregan en un pendrive provisto por el
  paciente, no por el centro.
\end{planos}
\end{CajaDestacada}

% =====================================================================
\section{Políticas y reglas del centro}
\label{sec:politicas}

\subsection{Política de cotización}

\begin{CajaDatos}
\boxtitle{Qué información necesita el centro para entregar un valor}
El centro necesita tres datos, en este orden de importancia:

\begin{center}
\begin{tabularx}{\textwidth}{@{}p{0.8cm}p{4.4cm}X@{}}
\toprule
\rowcolor{fondotabla}\textbf{N} & \textbf{Dato} &
\textbf{Para qué lo pide} \\
\midrule
1 & Qué examen se necesita & Determina el arancel base y si hace falta orden
  médica. \\
2 & Qué previsión tiene el paciente & Determina el valor. Sin este dato el
  centro puede negarse a dar un número. \\
3 & Si el paciente es adulto mayor & Ajusta la recomendación de acompañante y
  el bloque horario. \\
\bottomrule
\end{tabularx}
\end{center}
\end{CajaDatos}

\begin{CajaDatos}
\boxtitle{Cuándo necesita conocer la previsión}
El centro pide la previsión \textbf{siempre antes de dar el valor}, excepto
en dos casos:

\begin{hitos}
\item \textbf{Curvimetría:} el centro puede dar el valor particular sin
  preguntar, porque el arancel particular y el de Fonasa difieren en menos
  de \pesos{5.000} y el centro prefiere simplificar la llamada. Si el paciente
  menciona su previsión, el centro lo registra.
\item \textbf{Pacientes mayores de 70 años:} el centro pregunta la previsión
  antes que cualquier otra cosa, porque sabe que en la mayoría de los casos
  es Fonasa.
\end{hitos}
\end{CajaDatos}

\begin{CajaDatos}
\boxtitle{Cuándo puede entregar un valor particular}
El centro entrega el valor particular cuando el paciente dice que es
particular, que no tiene previsión o que paga de su bolsillo, y también
cuando el paciente no responde la pregunta sobre previsión. En ese uúltimo
caso entrega el valor particular y aclara: \emph{ese es el valor particular;
si tiene Fonasa o Isapre el valor cambia, me confirma y le doy el correcto}.
\end{CajaDatos}

\begin{CajaDatos}
\boxtitle{Qué ocurre si el paciente no indica su previsión}
El centro no insiste más de una vez. Si el paciente vuelve a no decir nada,
entrega el valor particular y marca la cotización como \textbf{particular
por omisión}. No vuelve a preguntar.
\end{CajaDatos}

\begin{CajaDatos}
\boxtitle{Qué información puede quedar pendiente de confirmación}
Cuatro datos pueden quedar como \textbf{no confirmado}:

\begin{center}
\begin{tabularx}{\textwidth}{@{}p{4.6cm}X@{}}
\toprule
\rowcolor{fondotabla}\textbf{Dato} & \textbf{Cuándo queda pendiente} \\
\midrule
Valor con Isapre & Cuando la isapre exige autorización previa por carta. El
  centro entrega un rango y la autorización queda pendiente. \\
Próxima hora & Cuando el centro pide al paciente que vuelva a llamar más
  tarde para confirmar. \\
Disponibilidad de sábado & Cuando la consulta se hizo un viernes después de
  las 16:00 y el sábado no se informa por teléfono. \\
Costo de cancelación & Cuando el paciente no aclara si es particular o con
  cobertura, y el valor depende de la modalidad. \\
\bottomrule
\end{tabularx}
\end{center}
\end{CajaDatos}

\subsection{Política de disponibilidad}
\label{sec:disp}

\begin{CajaDatos}
\boxtitle{Cómo se informa la disponibilidad}
El centro informa la disponibilidad por el mismo canal en que se hizo la
consulta. Entrega la fecha y el bloque, no una tabla completa de agenda.
\end{CajaDatos}

\begin{CajaDatos}
\boxtitle{Qué significa una hora disponible}
Una hora disponible es un bloque de agenda del examen solicitado, en un día
en que ese examen opera, dentro de los horarios de la sección
\ref{sec:horarios} y con la anticipación mínima cumplida. No es una reserva:
es una referencia que puede ser reasignada en cualquier momento.
\end{CajaDatos}

\begin{CajaDatos}
\boxtitle{Cuánto tiempo se mantiene una hora}
La hora se mantiene \textbf{48 horas} desde que se informa. Dentro de ese
plazo, si el paciente decide tomarla, debe confirmar por el canal que elija.
Pasadas las 48 horas, el centro la libera y puede entregarsela a otro
paciente, sin aviso.
\end{CajaDatos}

\begin{CajaDatos}
\boxtitle{Qué ocurre si una hora deja de estar disponible}
\begin{planos}
\item Si el paciente ya confirmó y el centro la libera, el centro se
  compromete a llamar para ofrecer una hora alternativa.
\item Si el paciente no había confirmado, la hora deja de estar disponible y
  el paciente debe volver a consultar.
\item El centro nunca avisa de forma proactiva que perdio una hora.
\end{planos}
\end{CajaDatos}

\begin{CajaDatos}
\boxtitle{Qué información entrega el centro al consultar disponibilidad}
Al consultar disponibilidad, el centro entrega:

\begin{hitos}
\item la fecha del bloque disponible;
\item la hora de inicio y el bloque (mañana o tarde);
\item la duración estimada del examen;
\item si el examen requiere orden médica;
\item el tiempo de anticipación que resta para confirmar.
\end{hitos}
Lo que no puede confirmar, lo declara como pendiente.
\end{CajaDatos}

\subsection{Política de reserva}

\begin{CajaNota}
Las reservas se confirman solo por teléfono o WhatsApp, en horario de
atención, y quedan cerradas cuando el centro entrega la confirmación al
paciente.
\end{CajaNota}

\begin{CajaDatos}
\boxtitle{Qué datos solicita el centro para reservar}
\begin{center}
\begin{tabularx}{\textwidth}{@{}p{4.6cm}X@{}}
\toprule
\rowcolor{fondotabla}\textbf{Dato} & \textbf{Para qué lo solicita} \\
\midrule
Nombre completo del paciente & Emitir la confirmación. \\
Previsión y número de ficha & Determinar el arancel final. \\
Orden médica & Validar que el examen se puede realizar. \\
Canal de contacto elegido & Enviar la confirmación. \\
\bottomrule
\end{tabularx}
\end{center}
\end{CajaDatos}

\begin{CajaDatos}
\boxtitle{En qué momento solicita el nombre del paciente}
El centro solicita el nombre \textbf{después} de acordar el examen, el valor
y la hora, en el momento en que el paciente decide tomar esa hora.
\end{CajaDatos}

\begin{CajaDatos}
\boxtitle{Qué datos de contacto solicita}
WhatsApp o teléfono, uno de los dos, para confirmar. Nunca correo y teléfono
a la vez. El RUT se solicita recien en recepción, el día del examen.
\end{CajaDatos}

\begin{CajaDatos}
\boxtitle{Qué información es necesaria para confirmar una hora}
\begin{hitos}
\item Examen y fecha u hora acordados.
\item Confirmación del paciente de que toma esa hora.
\item Previsión declarada.
\item Orden médica vigente, si el examen la exige.
\end{hitos}
Con esos cuatro elementos la reserva queda completa. Sin el primero no hay
reserva.
\end{CajaDatos}

\subsection{Política de datos personales}
\label{sec:datos}

\begin{CajaDestacada}
\boxtitle{Durante una simple cotización}
\begin{planos}
\item \textbf{No se requiere RUT.} El arancel depende del examen y de la
  modalidad, no de la identidad de la persona.
\item \textbf{No se requieren datos personales sensibles.} Nada de
  diagnósticos, recetas, antecedentes clínicos ni información de salud.
\item \textbf{No se solicita información innecesaria.} Nombre, domicilio y
  correo no son necesarios para saber si el centro hace el examen y cuanto
  cuesta.
\item \textbf{Los datos personales solo se solicitan en la etapa de
  reserva.}
\item \textbf{Los ejemplos de este documento no contienen datos de
  personas.}
\end{planos}
\end{CajaDestacada}

% =====================================================================
\section{Catálogo de exámenes}
\label{sec:catalogo}

% ---------------------------------------------------------- examen 1
\subsection{Curvimetría con optometría clínica}

\subsubsection{Identificación}

\begin{CajaDatos}
\begin{campos}
\item[Código] OFT-CV-001
\item[Nombre técnico] Curvimetría con optometría clínica y prueba de
  estereopsis.
\item[Nombre comun] Curvimetría.
\item[Formas coloquiales] curvimetría; examen de la vista; medir la vista;
  examen de optometría; chequeo de los ojos; examen para lentes; revisión de
  la graduación.
\item[Categoría] óptica y refracción.
\end{campos}
\end{CajaDatos}

\subsubsection{Descripción}

\begin{CajaDatos}
Es un examen no invasivo que mide el poder de refracción de ambos ojos y
evalua como trabajan juntos. \textbf{Para qué sirve:} para saber cuantas
dioptrías necesita la persona y si existen diferencias entre un ojo y otro.
\textbf{Permite obtener:} graduación por ojo, agudeza visual de lejos y de
cerca, visión en relieve y capacidad de fusión binocular. Se solicita de
forma habitual cuando el paciente quiere renovar lentes por primera vez o
cada dos años, cuando nota cambios de visión, o en los controles escolares.
\end{CajaDatos}

\subsubsection{Requisitos}

\begin{CajaDatos}
\begin{center}
\begin{tabularx}{\textwidth}{@{}p{5.2cm}X@{}}
\toprule
\rowcolor{fondotabla}\textbf{Requisito} & \textbf{Condición} \\
\midrule
Orden médica & \textbf{No requiere.} Se puede pedir sin ningún documento
  médico. \\
Ayuno & \textbf{No requiere.} \\
Preparación previa & \textbf{Ninguna.} No hay gotas, no hay dilatación, no
  hay que ajustar nada. \\
Actividades a suspender o modificar & \textbf{Ninguna.} Se puede venir
  trabajando, conduciendo o haciendo deporte. \\
Lentes de contacto & \textbf{Sin restricción.} Se puede venir con lentes de
  contacto puestos; el optometrista los retira durante la medición. \\
Acompañante & \textbf{No se requiere.} \\
Adultos mayores & \textbf{Sin requisito adicional.} Se ofrece silla con
  respaldo y se hacen pausas cada 10 minutos. \\
Otras condiciones & En menores de 8 años se requiere la presencia de un
  adulto responsable durante todo el examen. \\
\bottomrule
\end{tabularx}
\end{center}
\end{CajaDatos}

\subsubsection{Atención}

\begin{CajaDatos}
\begin{campos}
\item[Duración aproximada] 30 a 40 minutos.
\item[Modalidad] Presencial unicamente.
\item[Horario disponible] Lunes a viernes 08:30 a 13:00; martes, jueves y
  viernes 15:00 a 19:00; sábado 09:00 a 13:00.
\item[Días disponibles] Lunes a sábado (el sábado solo por la mañana).
\item[Anticipación requerida] 24 horas.
\item[Capacidad diaria] 14 atenciones por día.
\end{campos}
\end{CajaDatos}

\subsubsection{Precios}

\begin{CajaDatos}
\begin{center}
\begin{tabularx}{\textwidth}{@{}p{4.0cm}p{3.4cm}X@{}}
\toprule
\rowcolor{fondotabla}\textbf{Modalidad} & \textbf{Valor} &
\textbf{Condición} \\
\midrule
Particular & \pesos{25.000} & Valor de lista. \\
Fonasa & \pesos{12.500} & Sin orden médica. Con orden médica el valor baja a
  \pesos{9.800}. \\
Isapre & Sin valor cerrado & El centro no puede dar un valor sin ver la
  carta del contrato. Informa un rango de referencia entre \pesos{14.000} y
  \pesos{19.000} y pide autorización previa. \\
Convenio PSVM & \pesos{18.000} & Valor uúnico, sin autorización previa. \\
\bottomrule
\end{tabularx}
\end{center}
\end{CajaDatos}

\subsubsection{Resultados}

\begin{CajaDatos}
\begin{campos}
\item[Tipo de resultado] Informe de optometría con la graduación.
\item[Formato] Papel original firmado por el optometrista, más copia en PDF
  entregada en la memoria USB del paciente, si este trae una.
\item[Plazo de entrega] Inmediato: 30 minutos después del examen, en
  recepción.
\item[Incluye informe] Sí.
\item[Incluye imagenes] No.
\item[Recomendación posterior] Si la graduación cambia más de 1.00 dioptría
  respecto de la receta anterior, el optometrista recomienda renovar lentes y
  controlar en 6 meses.
\end{campos}
\end{CajaDatos}

\subsubsection{Recomendaciones}

\begin{CajaDatos}
\begin{planos}
\item \textbf{Antes:} llegar con la uúltima receta de lentes a la mano, si se
  tiene. No es obligatorio. Si no hay recetas previas, se registra que es la
  primera medición en el centro.
\item \textbf{Durante:} el examen no molesta. Se le pide al paciente que mire
  a una luz lejana y luego a una cercana. Puede parpadear libremente.
\item \textbf{Después:} no hay restricciones. Se puede volver al trabajo o
  seguir conduciendo normalmente.
\item \textbf{Consultar previamente:} antes del examen, si el paciente tiene
  glaucoma avanzado, cirugía ocular reciente (menos de 6 semanas),
  desprendimiento de retina previo, o una lesión corneal sin resolver.
\end{planos}
\end{CajaDatos}

\subsubsection{Conducción y acompañamiento}

\begin{CajaDatos}
\begin{center}
\begin{tabularx}{\textwidth}{@{}p{5.6cm}X@{}}
\toprule
\rowcolor{fondotabla}\textbf{Pregunta} & \textbf{Respuesta} \\
\midrule
Puede conducir después? & \textbf{Sí.} No hay dilatación, no hay sedación y
  no hay visión borrosa. \\
Se recomienda acompañante? & \textbf{No.} \\
Alguna condición que lo modifique? & \textbf{Si el paciente es menor de 14
  años} o es su primera medición, el centro sugiere que venga un adulto
  responsable, pero no lo exige. \\
\bottomrule
\end{tabularx}
\end{center}
\end{CajaDatos}

% ---------------------------------------------------------- examen 2
\subsection{Medición de glaucoma}

\subsubsection{Identificación}

\begin{CajaDatos}
\begin{campos}
\item[Código] OFT-GLU-002
\item[Nombre técnico] Tonometría no contactil (Goldmann) y campimetría visual
  computarizada.
\item[Nombre comun] Medición de glaucoma.
\item[Formas coloquiales] examen de glaucoma; medición de glaucoma; la
  presión del ojo; examen de la presión; campimetría; campo visual; examen de
  los nervios del ojo; revisión de glaucoma.
\item[Categoría] Neuro-oftalmología y glaucoma.
\end{campos}
\end{CajaDatos}

\subsubsection{Descripción}

\begin{CajaDatos}
Es un conjunto de dos mediciones. La primera, la tonometría, mide la presión
intraocular sin tocar el ojo, con un soplo de aire. La segunda, la
campimetría, mide el campo visual: que parte del campo ve el paciente y que
partes tiene perdidas, y lo representa en una rejilla. \textbf{Para que
sirve:} para detectar y seguir el glaucoma, que es la enfermedad que daña el
nervio óptico sin que el paciente note nada al principio. \textbf{Permite
obtener:} presión intraocular por ojo, umbral de sensibilidad, campimetría de
ambos ojos y un índice de daño del nervio óptico. Se solicita cuando hay
antecedentes familiares, cuando se sospecha glaucoma, en los controles de
pacientes ya diagnosticados, y a partir de los 40 años en personas con
factores de riesgo.
\end{CajaDatos}

\subsubsection{Requisitos}

\begin{CajaDatos}
\begin{center}
\begin{tabularx}{\textwidth}{@{}p{5.2cm}X@{}}
\toprule
\rowcolor{fondotabla}\textbf{Requisito} & \textbf{Condición} \\
\midrule
Orden médica & \textbf{Sí, siempre.} Orden vigente de un médico que indique
  medición de glaucoma o campimetría. Sin orden, el examen no se realiza bajo
  ningún supuesto. \\
Ayuno & \textbf{No requiere.} Se puede comer normalmente antes. \\
Preparación previa & \textbf{Sí, y es la más exigente del catálogo.}
  \begin{itemize}[label=\textbullet,leftmargin=1.1em,itemsep=1pt,topsep=2pt]
  \item Suspender lentes de contacto blandos 5 días antes.
  \item Suspender lentes rigidos o gas permeables 21 días antes.
  \item No usar gotas dilatadoras en las 48 horas previas.
  \item Traer una lista de los medicamentos que toma.
  \end{itemize} \\
Actividades a suspender o modificar & \textbf{No hay que suspender
  actividades.} Se puede venir trabajando. Solo se pide que la uúltima hora
  del examen se haga en silencio. \\
Lentes de contacto & \textbf{Prohibidos durante el examen} y durante los 5
  días previos. El paciente debe venir con anteojos. \\
Acompañante & \textbf{No obligatorio.} Se recomienda, por comodidad, si el
  paciente tiene 80 años o más, o si le cuesta mantener el silencio
  prolongado que exige la campimetría. \\
Adultos mayores & Se recomienda pedir hora en el primer bloque de la mañana
  (08:30 a 09:30), porque la campimetría exige mantener los ojos fijos
  durante varios minutos y la energía suele estar mejor en la mañana. El
  acompañante puede esperar en la sala de espera, no en la cámara. \\
Otras condiciones & Dificultad para distinguir bien de cerca o glaucoma en
  un ojo uúnico deben informarse al llegar. \\
\bottomrule
\end{tabularx}
\end{center}
\end{CajaDatos}

\subsubsection{Atención}

\begin{CajaDatos}
\begin{campos}
\item[Duración aproximada] 80 a 100 minutos (ambas mediciones, ambos ojos, en
  un mismo bloque).
\item[Modalidad] Presencial. Incluye una entrevista clínica de 10 minutos al
  inicio, que puede hacerse por telemedicina con 24 horas de anticipación.
\item[Horario disponible] Martes, jueves y viernes, 08:30 a 13:00.
\item[Días disponibles] \textbf{Solo martes, jueves y viernes.} No hay
  atenciones los lunes, miércoles, sabados ni domingos.
\item[Anticipación requerida] 5 días hábiles. Con menos anticipación, el
  centro no garantiza hora.
\item[Capacidad diaria] 6 pacientes por día, porque cada bloque de cámara
  ocupa 90 minutos.
\end{campos}
\end{CajaDatos}

\subsubsection{Precios}

\begin{CajaDatos}
\begin{center}
\begin{tabularx}{\textwidth}{@{}p{4.0cm}p{3.6cm}X@{}}
\toprule
\rowcolor{fondotabla}\textbf{Modalidad} & \textbf{Valor} &
\textbf{Condición} \\
\midrule
Particular & \pesos{72.000} & Valor de lista, sin autorización. \\
Fonasa & \pesos{45.000} & Con orden médica vigente, pagada en el acto. \\
Isapre & \pesos{58.000} referencial & \textbf{Solo referencial.} El valor
  depende de la carta de cobertura y de la autorización previa de la isapre,
  que puede demorar de 3 a 10 días hábiles. Si la isapre rechaza, se cobra el
  valor particular o Fonasa, según corresponda. \\
Convenio PSVM & \pesos{62.000} & Valor uúnico sin autorización. \\
\bottomrule
\end{tabularx}
\end{center}
\end{CajaDatos}

\subsubsection{Resultados}

\begin{CajaDatos}
\begin{campos}
\item[Tipo de resultado] Informe con tonometría, curvas de campimetría e
  índice de glaucoma.
\item[Formato] Informe en papel con curvas gráficas, más copia en PDF enviada
  al correo a petición del paciente.
\item[Plazo de entrega] 3 días hábiles desde la realización del examen.
\item[Incluye informe] Si, obligatorio.
\item[Incluye imagenes] Sí. Incluye curvas de campimetría en PDF y una captura
  del campo visual. No incluye fotografías del fondo del ojo: esas son de
  otro examen.
\item[Recomendación posterior] Si el índice de glaucoma supera el rango de
  referencia, el informe recomienda una consulta con oftalmólogo antes de 30
  días.
\end{campos}
\end{CajaDatos}

\subsubsection{Recomendaciones}

\begin{CajaDatos}
\begin{planos}
\item \textbf{Antes:} llegar 20 minutos antes (más que en los otros exámenes,
  porque hay entrevista). Traer la lista de medicamentos y los anteojos.
  Haber comido normalmente. No suspender las gotas habituales, salvo las
  dilatadoras.
\item \textbf{Durante:} el examen es indoloro. La campimetría pide mantener
  la mirada fija en un punto central mientras se registra automáticamente lo
  que ve el paciente; el paciente no debe mover la cabeza ni parpadear de
  forma agresiva. Si se marea, puede pedir una pausa de 5 minutos sin costo.
\item \textbf{Después:} puede hacer todo lo que quiera. No hay visión
  borrosa ni dilatación.
\item \textbf{Consultar previamente:} si el paciente tiene glaucoma ya
  diagnosticado y está recién iniciando tratamiento, si fue operado del ojo
  en los uúltimos 30 días, si tiene una inflamación activa del ojo, o si toma
  corticoides orales de forma crónica.
\end{planos}
\end{CajaDatos}

\subsubsection{Conducción y acompañamiento}

\begin{CajaDatos}
\begin{center}
\begin{tabularx}{\textwidth}{@{}p{5.6cm}X@{}}
\toprule
\rowcolor{fondotabla}\textbf{Pregunta} & \textbf{Respuesta} \\
\midrule
Puede conducir después? & \textbf{Sí.} Este examen no utiliza dilatación ni
  sedación. \\
Se recomienda acompañante? & \textbf{Sugerido, no obligatorio}, para pacientes
  de 80 años o más, o con dificultad para el silencio prolongado de la
  campimetría. \\
Alguna condición que lo modifique? & Ninguna. La recomendación de acompañante
  no cambia según la previsión ni según el resultado. \\
\bottomrule
\end{tabularx}
\end{center}
\end{CajaDatos}

% ---------------------------------------------------------- examen 3
\subsection{Fondo de ojo}

\subsubsection{Identificación}

\begin{CajaDatos}
\begin{campos}
\item[Código] OFT-FDO-003
\item[Nombre técnico] Fundoscopia con oftalmoscopio indirecto y dilatación
  pupilar.
\item[Nombre comun] Fondo de ojo.
\item[Formas coloquiales] fondo de ojo; examen del fondo de ojo; examen de la
  retina; revisión de la retina; mirar la retina; examen del nervio óptico;
  foto del fondo del ojo.
\item[Categoría] Retina y fondo de ojo.
\end{campos}
\end{CajaDatos}

\subsubsection{Descripción}

\begin{CajaDatos}
Es el examen en el que se mira el interior del ojo con un oftalmoscopio
indirecto, previa dilatación de la pupila. \textbf{Para qué sirve:} para ver
la retina, el disco del nervio óptico y la córnea, y detectar retinopatías,
degeneraciones, hemorragias o signos de hipertensión. \textbf{Permite
obtener:} un mapa del fondo de ojo, el estado del nervio óptico y de la
retina, y una fotografía digitalizada de referencia. Se solicita en los
controles de diabetes (incluida la diabetes gestacional), en los controles de
hipertensión, en pacientes que usan corticoides de forma crónica, o cuando
hay una pérdida de visión que no se explica por la graduación.
\end{CajaDatos}

\subsubsection{Requisitos}

\begin{CajaDatos}
\begin{center}
\begin{tabularx}{\textwidth}{@{}p{5.2cm}X@{}}
\toprule
\rowcolor{fondotabla}\textbf{Requisito} & \textbf{Condición} \\
\midrule
Orden médica & \textbf{Depende de la modalidad de pago.} Con Fonasa o
  convenio: obligatoria. Con Isapre: obligatoria y con autorización previa.
  \textbf{A particular: no se exige.} \\
Ayuno & \textbf{Sí exige ayuno de 6 horas.} En ayunas se evita que el reflejo
  del estómago altere la medida asociada al examen. \\
Preparación previa & \textbf{Sí.}
  \begin{itemize}[label=\textbullet,leftmargin=1.1em,itemsep=1pt,topsep=2pt]
  \item Ayuno de 6 horas, solo agua.
  \item Suspender las gotas que contengan atropina o fenilefrina al menos
    48 horas antes, solo si fueron indicadas por un especialista.
  \item Informar con anticipación si el paciente tiene glaucoma y si hay
    alergia al fármaco dilatador que usa el centro.
  \end{itemize} \\
Actividades a suspender o modificar & \textbf{Sí.} El paciente no debe
  conducir, no debe subir escaleras sin ayuda y no debe operar maquinaria
  durante las 4 a 6 horas posteriores, porque la pupila queda dilatada. \\
Lentes de contacto & \textbf{Se deben suspender 24 horas antes.} Durante el
  examen y después de el hay que usar anteojos. \\
Acompañante & \textbf{Sí, recomendado de forma expresa.} No por el examen en
  si, sino por las 4 a 6 horas de visión borrosa posteriores. Si el paciente
  viene solo, el centro le pide que vuelva a casa en un medio de transporte
  público o en un taxi. \\
Adultos mayores & \textbf{Acompañante obligatorio en la práctica.} Por el
  riesgo de caidas al subir o bajar escalones con visión borrosa y por el
  riesgo de desorientación. Si no hay acompañante, el centro prioriza la
  primera hora de la mañana. \\
Otras condiciones & En caso de glaucoma de anángulo cerrado, el centro exige
  una carta del oftalmólogo tratante antes de dilatar la pupila. \\
\bottomrule
\end{tabularx}
\end{center}
\end{CajaDatos}

\subsubsection{Atención}

\begin{CajaDatos}
\begin{campos}
\item[Duración aproximada] 25 a 30 minutos en la cámara, más 30 a 45 minutos
  de visión borrosa posterior que el paciente pasa sentado en recepción.
\item[Modalidad] Presencial.
\item[Horario disponible] Lunes a viernes 08:30 a 13:00; además miércoles
  15:00 a 19:00.
\item[Días disponibles] Lunes a viernes, más el miércoles por la tarde.
\item[Anticipación requerida] 3 días hábiles.
\item[Capacidad diaria] 12 atenciones por día, pero solo 6 después de las
  15:00 (solo el miércoles).
\end{campos}
\end{CajaDatos}

\subsubsection{Precios}

\begin{CajaDatos}
\begin{center}
\begin{tabularx}{\textwidth}{@{}p{4.0cm}p{3.4cm}X@{}}
\toprule
\rowcolor{fondotabla}\textbf{Modalidad} & \textbf{Valor} &
\textbf{Condición} \\
\midrule
Particular & \pesos{34.000} & Sin necesidad de orden médica. \\
Fonasa & \pesos{19.500} & \textbf{Requiere orden médica.} Sin orden médica se
  cobra el valor particular de \pesos{34.000}. \\
Isapre & \pesos{21.000} & Requiere orden médica y autorización previa por
  carta. La autorización puede demorar de 3 a 10 días hábiles. \\
Convenio PSVM & No cubierto & El convenio no incluye fondo de ojo. Si el
  paciente quiere usarlo, se le cobra el valor particular de \pesos{34.000},
  y se le informa antes de agendar. \\
\bottomrule
\end{tabularx}
\end{center}
\end{CajaDatos}

\subsubsection{Resultados}

\begin{CajaDatos}
\begin{campos}
\item[Tipo de resultado] Informe de fondo de ojo con descripción de retina,
  nervio óptico y córnea, más una fotografía digitalizada.
\item[Formato] Informe en papel firmado por el oftalmólogo, una impresión en
  color de la fotografía, y el archivo JPG en un pendrive del paciente.
\item[Plazo de entrega] 2 días hábiles. La fotografía se entrega el mismo
  día; el informe escrito, dos días después.
\item[Incluye informe] Sí.
\item[Incluye imagenes] Sí. Una fotografía por ojo, más una por sector si el
  oftalmólogo lo indica.
\item[Recomendación posterior] Si el informe detecta una lesión sospechosa, la
  recomendación es control con oftalmólogo en un plazo de 30 días, no antes.
\end{campos}
\end{CajaDatos}

\subsubsection{Recomendaciones}

\begin{CajaDatos}
\begin{planos}
\item \textbf{Antes:} ayunar 6 horas. Traer anteojos, porque el paciente va a
  salir con visión borrosa y no puede usar lentes de contacto. Venir
  acompañado. Si usa lentes de contacto, dejarlos en casa desde la noche
  anterior. Informar si tiene glaucoma o alergias farmacéuticas.
\item \textbf{Durante:} se instilan gotas dilatadoras en cada ojo, lo que
  produce un ardor breve y hace ver borroso. El examen es indoloro. El
  paciente ve una luz brillante y no puede enfocar durante 10 minutos. Puede
  parpadear normalmente.
\item \textbf{Después:} 4 a 6 horas de visión borrosa y sensibilidad a la
  luz. El centro sugiere anteojos de sol. Conviene permanecer en el centro
  unos 30 minutos antes de salir. Puede comer y beber con normalidad.
\item \textbf{Consultar previamente:} si el paciente tiene glaucoma de anángulo
  cerrado, si fue operado de cataratas en los uúltimos 30 días, si presenta un
  ojo muy rojo o con dolor, si tiene retinopatía diabética severa
  diagnosticada, o si tiene menos de 6 meses de cirugía ocular.
\end{planos}
\end{CajaDatos}

\subsubsection{Conducción y acompañamiento}

\begin{CajaDatos}
\begin{center}
\begin{tabularx}{\textwidth}{@{}p{5.6cm}X@{}}
\toprule
\rowcolor{fondotabla}\textbf{Pregunta} & \textbf{Respuesta} \\
\midrule
Puede conducir después? & \textbf{No puede conducir} al menos 4 horas, y
  hasta 6 si la pupila responde lento. \\
Se recomienda acompañante? & \textbf{Sí, de forma expresa y por defecto.} La
  recomendación aplica a todos los pacientes, sin importar la edad ni la
  previsión. \\
Alguna condición que lo modifique? & La recomendación es más estricta en
  pacientes mayores de 70 años, en diabéticos y en personas con problemas de
  equilibrio: en esos casos el acompañante se declara obligatorio y no se
  acepta al paciente sin el. \\
\bottomrule
\end{tabularx}
\end{center}
\end{CajaDatos}

% ------------------------------------------------- resumen comparativo
\subsection{Resumen comparativo}

\begin{CajaDatos}
\begin{center}
\small
\begin{tabularx}{\textwidth}{@{}p{3.7cm}X X X@{}}
\toprule
\rowcolor{fondotabla}\textbf{Criterio} & \textbf{OFT-CV-001} &
\textbf{OFT-GLU-002} & \textbf{OFT-FDO-003} \\
\midrule
Examen & Curvimetría & Medición de glaucoma & Fondo de ojo \\
Orden médica & No & Siempre & Según modalidad \\
Ayuno & No & No & Sí, 6 horas \\
Preparación previa & Ninguna & 5 a 21 días & 24 h más ayuno \\
Lentes de contacto & Sin restricción & Prohibidos 5 a 21 días &
  Suspender 24 h \\
Acompañante & No & Sugerido si 80 o más & Obligatorio \\
Puede conducir & Sí & Sí & No, 4 a 6 h \\
Duración & 30 a 40 min & 80 a 100 min & 25 a 30 min \\
Días & L a sábado & Mar, jue y vie & L a vie, mie tarde \\
Anticipación & 24 h & 5 días hábiles & 3 días hábiles \\
Capacidad diaria & 14 & 6 & 12 (6 por tarde) \\
Particular & \pesos{25.000} & \pesos{72.000} & \pesos{34.000} \\
Fonasa & \pesos{12.500} & \pesos{45.000} & \pesos{19.500} \\
Isapre & Rango sin valor cerrado & Referencial \pesos{58.000} &
  \pesos{21.000} con autorización \\
Convenio PSVM & \pesos{18.000} & \pesos{62.000} & No cubierto \\
Plazo de resultado & Inmediato & 3 días hábiles & 2 días hábiles \\
Incluye imágenes & No & Sí, curvas & Sí, fotografías \\
Informe & Sí & Sí, obligatorio & Sí \\
\bottomrule
\end{tabularx}
\end{center}
\end{CajaDatos}

\end{document}
